\chapter{Euler--Bernoulli Beam Equation}

\section*{Introduction}

When you stand on a diving board, it bends. The deflection depends on your weight (the load), the board's length, and its stiffness (resistance to bending). The Euler--Bernoulli equation tells you exactly how much the board deflects at every point and how it vibrates after you jump off. The key insight is that bending involves a balance between internal bending moment (resisting deformation) and external load (causing it), with inertia (mass $\times$ acceleration) opposing the motion during vibration.
\section*{Fundamental Equation Used}

For a beam of flexural rigidity $EI$, cross-sectional area $A$, and density $\rho$, the transverse deflection $w(x,t)$ satisfies:
\[
EI\frac{\partial^4 w}{\partial x^4} + \rho A \frac{\partial^2 w}{\partial t^2} = q(x,t)
\]
\section*{Assumptions}

\textbf{Assumptions:} The beam is slender (length $\gg$ depth and width). Cross-sections remain plane and perpendicular to the neutral axis (Euler--Bernoulli hypothesis). Shear deformation and rotary inertia are neglected. The material is linear elastic (stress $\propto$ strain).


\textbf{Limitations:} Inaccurate for short, deep beams ($L/h < 10$) where shear deformation matters (use Timoshenko beam theory). Does not account for geometric nonlinearity (large deflections), material nonlinearity, or buckling. Fails at very high frequencies where rotatory inertia matters.


\section*{Mathematical Derivation}

\textbf{Step 1: Define the beam geometry and the Euler--Bernoulli hypothesis.}

Consider a slender beam of length $L$ with cross-sectional area $A$, Young's modulus $E$, and second moment of area $I = \int y^2\,dA$ about the neutral axis. Let $w(x,t)$ be the transverse deflection at position $x$ along the beam and time $t$. The Euler--Bernoulli hypothesis states: cross-sections remain plane and perpendicular to the deformed neutral axis (shear deformation is neglected).

\textbf{Step 2: Relate curvature to bending moment.}

For small deflections, the curvature of the neutral axis is $\kappa \approx d^2w/dx^2$. The bending moment at a cross-section is proportional to the curvature:
\begin{equation}
M(x,t) = EI\frac{\partial^2 w}{\partial x^2}
\end{equation}
where $EI$ is the flexural rigidity (a measure of the beam's resistance to bending). This is the Euler--Bernoulli bending law, valid for linear elastic materials.

\textbf{Step 3: Relate shear force to load.}

Consider equilibrium of a beam element of length $\delta x$. The shear force $V$ satisfies:
\begin{equation}
\frac{\partial V}{\partial x} = q(x,t) - \rho A\frac{\partial^2 w}{\partial t^2}
\end{equation}
where $q(x,t)$ is the distributed transverse load per unit length and $-\rho A\,\partial^2 w/\partial t^2$ is the inertial force per unit length (the negative sign arises because the acceleration is in the direction of $w$, while the reaction force opposes it).

\textbf{Step 4: Relate shear force to bending moment.}

Moment equilibrium of the beam element (taking moments about one end and neglecting rotary inertia) gives:
\begin{equation}
V = \frac{\partial M}{\partial x}
\end{equation}

\textbf{Step 5: Combine the relationships.}

Differentiating $V = \partial M/\partial x$:
\begin{equation}
\frac{\partial V}{\partial x} = \frac{\partial^2 M}{\partial x^2} = \frac{\partial^2}{\partial x^2}\left(EI\frac{\partial^2 w}{\partial x^2}\right) = EI\frac{\partial^4 w}{\partial x^4}
\end{equation}
where we assumed $EI$ is constant along the beam. Equating with Step 3:
\begin{equation}
EI\frac{\partial^4 w}{\partial x^4} = q(x,t) - \rho A\frac{\partial^2 w}{\partial t^2}
\end{equation}

\textbf{Step 6: Write the final equation.}

Rearranging:
\begin{equation}
\boxed{EI\frac{\partial^4 w}{\partial x^4} + \rho A\frac{\partial^2 w}{\partial t^2} = q(x,t)}
\end{equation}
This is the Euler--Bernoulli beam equation: a fourth-order PDE in $x$ and second-order in $t$. The fourth spatial order reflects the chain: load $\to$ shear $\to$ moment $\to$ curvature $\to$ deflection (each step adds one derivative). Boundary conditions at each end require specifying two quantities (e.g., deflection and slope, or deflection and moment), totaling four conditions matching the fourth-order spatial derivative.

\section*{Characteristics of the Equation}

\begin{itemize}
\item \textbf{Fourth-order PDE:} The $x$-derivative is fourth order, requiring four boundary conditions (two at each end).
\item \textbf{Static deflection:} For time-independent loads: $EI \, d^4w/dx^4 = q(x)$, a fourth-order ODE.
\item \textbf{Free vibration:} Setting $q = 0$ and $w(x,t) = W(x)e^{i\omega t}$ gives $EI \, W'''' = \rho A \omega^2 W$, an eigenvalue problem for natural frequencies $\omega_n$ and mode shapes $W_n(x)$.
\item \textbf{Mode shapes:} For a simply supported beam: $W_n(x) = \sin(n\pi x/L)$, with $\omega_n = (n\pi)^2 \sqrt{EI/(\rho A L^4)}$.
\item \textbf{Superposition:} For linear elastic response, the total deflection is a sum of mode contributions.
\end{itemize}
\section*{Solution Methods}

\begin{enumerate}
\item \textbf{Direct integration:} For static loads, integrate $EI \, w'''' = q(x)$ four times, applying boundary conditions at each step.
\item \textbf{Eigenvalue expansion:} For vibration, expand $w(x,t) = \sum_n a_n W_n(x) e^{i\omega_n t}$ where $W_n$ are mode shapes satisfying boundary conditions.
\item \textbf{Rayleigh--Ritz method:} Assume a trial function satisfying boundary conditions, minimize the energy functional to approximate natural frequencies.
\item \textbf{Finite element:} Discretize the beam into elements, assemble stiffness and mass matrices, solve the generalized eigenvalue problem.
\end{enumerate}
\section*{Verifications}

\begin{itemize}
\item \textbf{Static deflection:} Measured deflections of simply supported beams under point loads match $PL^3/(48EI)$.
\item \textbf{Natural frequencies:} Measured frequencies of cantilever beams match $\omega_n = (\beta_n L)^2 \sqrt{EI/(\rho A L^4)}$ where $\beta_n L$ are the characteristic values.
\item \textbf{Mode shapes:} Laser vibrometry measurements of vibrating beams confirm the predicted sinusoidal (simply supported) or hyperbolic-trigonometric (cantilever) mode shapes.
\item \textbf{Resonance avoidance:} Bridge and building designs verified by full-scale vibration testing.
\end{itemize}
\section*{Applications}

\begin{itemize}
\item \textbf{Bridge design:} Deflection under traffic loads, natural frequencies to avoid resonance with wind or earthquake excitation.
\item \textbf{Building floors:} Vibration criteria for human comfort (floor frequencies above $\sim$8\,Hz to avoid perception).
\item \textbf{Aircraft wings:} Flutter analysis---coupling between aerodynamic forces and structural vibration.
\item \textbf{MEMS:} Micro-cantilever sensors detect mass changes (attogram sensitivity) through frequency shifts.
\item \textbf{Guitar strings:} Vibration modes determine the harmonic content of the sound.
\end{itemize}
\section*{What Richard Feynman Would Have Asked}

\subsection*{1. Plain-English Translation}
\textit{``What is the Euler--Bernoulli beam equation actually telling us?''}

It says: ``When you push on a beam, it bends in a way that balances the internal stiffness against the external load, with the beam's mass opposing the motion during vibration.'' The fourth derivative on the left side represents the beam's internal resistance to bending---how strongly the material fights back against deformation. The second time derivative represents inertia---the beam's resistance to acceleration. The right side is the external push (load).

Think of a diving board: when you stand on the end, the board bends because your weight (right side) exceeds the internal stiffness (left side) at that point. When you jump off, the board vibrates because inertia carries it past the equilibrium position, and stiffness pulls it back. The equation tracks this dance between stiffness, inertia, and external forcing.

\subsection*{2. Localized Mechanics}
\textit{``What is happening at a single cross-section of the beam as it bends?''}

At any cross-section, the upper fibers of the beam are in compression (shortened) and the lower fibers are in tension (elongated), or vice versa. The neutral axis (usually the centroid) has zero strain. The bending moment $M = EIw''$ measures how much the beam curves at that point---higher curvature means higher internal stress. The shear force $V = EIw'''$ measures how much the cross-section slides relative to its neighbors.

He would press you on why the equation is fourth-order: the physical chain is load $\to$ shear $\to$ moment $\to$ curvature $\to$ deflection. Each step involves a derivative: load is the derivative of shear, shear of moment, moment of curvature, and curvature is the second derivative of deflection. The fourth-order equation simply chains these relationships together.

\subsection*{3. Testing Extreme Limits}
\textit{``What happens when we push the beam to extreme conditions?''}

At very high frequencies (short wavelengths), the beam's response is dominated by shear deformation and rotatory inertia---the Euler--Bernoulli theory overestimates frequencies. The Timoshenko theory corrects this.

For very long, thin beams (like fishing rods), large deflections make the linear theory inaccurate---geometric nonlinearity (stretching of the neutral axis) must be included.

At the microscopic scale (MEMS cantilevers), the same equation applies, but the beam dimensions are micrometers and the loads are nanonewtons. The equation's universality means it works from meter-scale bridges to micrometer-scale sensors.

\subsection*{4. Dimensionality and Geometry}
\textit{``How does the cross-section shape affect the beam's behavior?''}

The flexural rigidity $EI$ depends on the cross-sectional shape through the second moment of area $I$. A hollow tube has higher $I$ than a solid rod of the same mass---this is why bicycle frames are hollow. An I-beam concentrates material far from the neutral axis, maximizing $I$ for minimum weight.

The orientation matters: a beam bent about its strong axis (high $I$) is much stiffer than when bent about its weak axis (low $I$). The equation handles this naturally through the appropriate $I$ value for the bending direction.

\subsection*{5. Cross-Disciplinary Connection}
\textit{``Where else does this same mathematical structure appear?''}

The Euler--Bernoulli equation is a special case of the more general plate equation ($D\nabla^4 w + \rho h \ddot{w} = q$) and shell equations. The same fourth-order structure appears in quantum mechanics (the biharmonic operator appears in certain potential problems), in fluid mechanics (the Stokes stream function for axisymmetric flow satisfies $\nabla^4 \psi = 0$), and in elasticity (the Airy stress function satisfies $\nabla^4 \phi = 0$ for 2D problems).

The deep connection is that any system where the restoring force depends on the curvature of the deformation (rather than the displacement itself) produces a fourth-order equation. This is because curvature involves second derivatives, and the equilibrium of bending moments requires two more derivatives.


%=============================================================
\section*{FAQ}

\textbf{Q1: Why is the equation fourth-order in $x$?}

A: Because bending involves the curvature ($w''$), and the bending moment is proportional to curvature ($M = EIw''$). The shear force is the derivative of the moment ($V = M' = EIw'''$), and the load is the derivative of the shear ($q = V' = EIw''''$). Each differentiation adds one order: curvature $\to$ moment $\to$ shear $\to$ load.

\textbf{Q2: What is the difference between simply supported and clamped?}

A: Simply supported: the beam can rotate at the supports but not translate ($w = 0$, $M = 0$). Clamped (fixed): neither translation nor rotation is allowed ($w = 0$, $w' = 0$). These different boundary conditions produce different mode shapes and natural frequencies. A clamped beam is stiffer and has higher frequencies than a simply supported beam of the same length.

\textbf{Q3: When does the Euler--Bernoulli theory fail?}

A: For short, deep beams ($L/h < 10$), shear deformation becomes significant and the Timoshenko beam theory (which includes shear and rotatory inertia) is needed. For large deflections (where $w$ is comparable to the beam length), geometric nonlinearity matters. At very high frequencies, the Timoshenko correction is also needed.
\section*{Important Bibliography}

\begin{enumerate}
\item Timoshenko, S.P., \textit{Theory of Plates and Shells} (McGraw-Hill, 1940), Chapter 1.
\item Weaver, W. and Gere, J.M., \textit{Matrix Analysis of Framed Structures}, 3rd ed. (Springer, 1990), Chapter 3.
\item Rao, S.S., \textit{Mechanical Vibrations}, 6th ed. (Pearson, 2017), Chapter 5.
\item Blevins, R.D., \textit{Formulas for Natural Frequency and Mode Shape} (Krieger, 2001), Chapter 2.
\end{enumerate}


